\section{Conditional synthesis}
\label{sec:conditional}

The following implication is the only global theorem of TPC-MVP1.
Its proof is short because the difficulty has been made explicit in
the hypotheses.

\begin{theorem}[\Ctag\ Conditional fixed-shift Hardy--Littlewood
synthesis]
\label{thm:conditional-hl}
Fix an admissible nonzero even integer \(h_0\) and
\(W\in C_c^\infty((0,\infty))\).  Assume H1--H9 for the literal TPC
carrier, with every estimate taken in the same physical
normalization and quantifier scope.  Then
\begin{equation}
\sum_{n\ge1}\Lambda(n)\Lambda(n+h_0)W(n/X)
=
\Ssing(h_0)X\int_0^\infty W(t)\,dt+o(X).
\label{eq:conditional-hl}
\end{equation}
\end{theorem}

\begin{proof}
Hypotheses H1, H4, and H6 produce an exhaustive subpower-size
classification of physical blocks.  Every term assigned to the soft
remainder contributes \(o(X)\) after the complete outer summation.
The two established exceptional returns together with H2 handle the
resonant blocks; H3 handles the generic signed blocks; and H4 returns
every remaining short, high, ultra, boundary, and outer block.  Thus
every nonsoft raw block entering synthesis satisfies
\[
\max_{b\in\BX^{\mathrm{ns}}}|\mathcal P_b|
\le X^{1-\sigma_*+o(1)}.
\]
Hypothesis H5 certifies that the determinant/zero-mode transfer used
in that statement has no unrecorded incompatibility.  H7 returns all
auxiliary parameters to the prescribed \(h_0\), and H9 gives
\[
\left|
\mathscr S_X\bigl(\{\mathcal P_b\}_{b\in\BX^{\mathrm{ns}}}\bigr)
\right|
\le
X^{1-\sigma_*+\Lambda_{\mathrm{phys}}+o(1)}
=o(X),
\]
because \(\Lambda_{\mathrm{phys}}<\sigma_*\).  H8 identifies this
synthesis output with \(B_{h_0,\delta}(X)\), up to the already
recorded \(o(X)\) remainder.  Therefore
\[
B_{h_0,\delta}(X)=o(X).
\]
Insert this in the established identity
\eqref{eq:hard-packet-anchor}.
\end{proof}

\begin{remark}[What the theorem does and does not prove]
\cref{thm:conditional-hl} proves an implication from explicit
unproved hypotheses to the target correlation.  It does not prove
H1--H9.  In particular H3 contains the unresolved growing
fixed-\(h_0\) arithmetic cancellation, and the conjunction of the
hypotheses may be stronger than \eqref{eq:conditional-hl}.
\end{remark}

\subsection{The conditional prime-pair consequence}

\begin{corollary}[\Ctag\ Conditional fixed-gap prime pairs]
\label{cor:conditional-pairs}
Assume H1--H9 for a fixed admissible even \(h_0\) and for one fixed
nonnegative, nonzero \(W\in C_c^\infty((1,2))\).  Then
\[
\#\{p\in[X,2X+O_{h_0}(1)]:
      p\ \text{and}\ p+h_0\ \text{are prime}\}
\gg_{h_0,W}\frac{X}{(\log X)^2}
\]
for all sufficiently large \(X\).
\end{corollary}

\begin{proof}
Since \(\Ssing(h_0)>0\) and \(\int W>0\),
\cref{thm:conditional-hl} gives
\[
\mathcal C_{h_0,W}(X)\gg_{h_0,W}X.
\]
The terms for which \(n\) or \(n+h_0\) is a prime power of exponent
at least two contribute at most
\[
O_{h_0,W}\!\left(X^{1/2}(\log X)^2\right)=o(X).
\]
The remaining mass comes from prime-prime pairs, and each term is at
most \(O_W((\log X)^2)\).  Division gives the stated lower bound.
\end{proof}

\begin{corollary}[\Ctag\ Conditional \(h_0=2\) specialization]
\label{cor:conditional-twins}
If H1--H9 are eventually proved for the prescribed shift \(h_0=2\),
then \cref{cor:conditional-pairs} implies infinitely many twin primes.
\end{corollary}

\begin{remark}[No sharp unsmoothing is needed for infinitude]
A single fixed nonnegative smooth weight already yields the
conditional lower bound in \cref{cor:conditional-pairs}.  Smooth
majorants and minorants may recover sharper interval formulations,
but positivity and removal of higher prime powers are not the new
parity-facing input.  That input is H3, or an equally strong signed
replacement for a failed positive route.
\end{remark}

\subsection{A proof skeleton in nine moves}

The conditional theorem may be read as the following end-to-end
program:
\begin{enumerate}[leftmargin=2em]
\item calibrate the singular-series main term and isolate
  \(B_{h_0,\delta}\);
\item export the literal coefficient into canonical native packets;
\item return the two established exceptional resonance branches;
\item prove or replace the \(W_X,\PX,\XX\) positive resonance gates;
\item prove signed generic affine M\"obius cancellation on the actual
  carrier;
\item return every short, high, ultra, boundary, and outer sector;
\item verify determinant energy and the distinguished zero mode;
\item reassemble and localize to the prescribed \(h_0\) with strict
  endpoint slack; and
\item reconnect to the original hard packet and use positivity only
  after the weighted asymptotic is obtained.
\end{enumerate}
No move may be deleted by renaming its output.
